\documentclass[a4paper,border=10mm]{standalone}
\usepackage[spanish, es-tabla]{babel}
\usepackage{tikz}
\usetikzlibrary{positioning,calendar,decorations.markings}

% --- Colores ---
\colorlet{darkgreen}{green!50!black}

% --- Traducción y Estilo de Cabeceras ---
\newcommand{\calrow}[1]{%
  \node[anchor=base east](Mon){L};
  \node[base right=of Mon](Tue){M}; \node[base right=of Tue](Wed){M};
  \node[base right=of Wed](Thu){J}; \node[base right=of Thu](Fri){V};
  \node[base right=of Fri](Sat){S}; \node[base right=of Sat](Sun){D};
  \node[darkgreen,above=of Thu]{\textbf{#1}};
}

\newcommand{\holidays}{% Feriados Argentina
  if (equals=01-01) [red!80!black] if (equals=03-24) [red!80!black]
  if (equals=04-02) [red!80!black] if (equals=05-01) [red!80!black]
  if (equals=05-25) [red!80!black] if (equals=06-20) [red!80!black]
  if (equals=07-09) [red!80!black] if (equals=12-08) [red!80!black]
  if (equals=12-25) [red!80!black]
}

% --- 1. Función para meses NORMALES ---
\newcommand{\calperiod}[1]{%
  \calendar[dates=\the\year-#1-01 to \the\year-#1-last]
  if (Sunday) [red!80!black]
  \holidays;
}

% --- 2. Función para meses CON RANGO (Sin líneas en blanco para
% evitar errores) ---
\newcommand{\calrange}[3]{%
  \calendar[dates=\the\year-#1-01 to \the\year-#1-last]
  if (Sunday) [red!80!black]
  \holidays
  if (between=\the\year-#1-#2 and \the\year-#1-#3)
  [days={fill=green!50, rectangle, inner sep=3pt}]
  if (equals=\the\year-#1-#2) [days={fill=red!50, rectangle, inner sep=3pt}]
  if (equals=\the\year-#1-#3) [days={fill=red!50, rectangle, inner sep=3pt}];
}

\begin{document}
\begin{tikzpicture}[every calendar/.style={week list}, scale=0.9]
  \sffamily
  \matrix[%
    row 1/.style={darkgreen,node distance=.3ex},
    row 3/.style={darkgreen,node distance=.3ex},
    row 5/.style={darkgreen,node distance=.3ex},
    column sep=2ex,
    row sep=3ex,
    draw=darkgreen,thick,rounded corners=5pt,
    postaction={decorate,decoration={markings,mark=at position 0.51 with
        {\node[fill=white,text=darkgreen,font={\bfseries\huge},anchor=west]
          (year)
    {\the\year{} Escuela 4-239 Cerro el Sosneado};}}}
  ]{%
    % Ejemplo: Usamos \calperiod para los normales y \calrange para
    % marcar rangos
    \calrow{Enero} & \calrow{Febrero} & \calrow{Marzo} &
    \calrow{Abril} \\
    \calperiod{01} & \calperiod{02} & \calrange{03}{16}{31} &
    \calrange{04}{17}{25} \\[1ex]

    \calrow{Mayo} & \calrow{Junio} & \calrow{Julio} & \calrow{Agosto} \\
    \calrange{05}{16}{31} & \calrange{06}{15}{30} & \calrange{07}{23}{31} &
    \calrange{8}{16}{31} \\[1ex]

    \calrow{Septiembre} & \calrow{Octubre} & \calrow{Noviembre} &
    \calrow{Diciembre} \\
    \calrange{09}{15}{30} & \calrange{10}{16}{24} &
    \calrange{11}{15}{30} & \calrange{12}{11}{19} \\
  };
\end{tikzpicture}
\end{document}
